The three classical apportionment paradoxes --- Alabama, Population,
and New-States --- are not pathological curiosities but inevitable
consequences of Hamilton's quota-compliance mechanism.
The non-monotone fractional-part ranking that enables Hamilton to
satisfy the quota rule for all possible population vectors is
precisely the mechanism that creates the paradoxes.

Divisor methods avoid this trap by using a different mechanism:
monotone priority queues that only add seats, never reconsidering
previously allocated seats.
The formal proof of divisor method immunity
(Theorem~\ref{thm:divisor-immune}) establishes that this immunity
holds for all four divisor methods (HH, Webster, Adams, Jefferson)
and for all possible population vectors.

The Balinski-Young impossibility theorem
(Theorem~\ref{thm:by-impossibility}) closes the loop:
paradox immunity and quota compliance are mathematically incompatible,
and no future mathematical innovation can resolve this incompatibility.
The tradeoff is real and irreducible.

For \textsc{bisect} platform users:
\begin{itemize}
  \item Huntington-Hill is provably paradox-free; any redistricting
    pipeline built on \textsc{bisect-apportion}'s HH implementation
    inherits this guarantee.
  \item The \texttt{check\_alabama\_paradox} function provides
    empirical verification for any method, enabling practitioners
    to audit non-HH apportionments.
  \item The Balinski-Young theorem should be cited whenever a party
    in redistricting litigation argues that a ``better'' apportionment
    method exists that is both quota-compliant and paradox-free ---
    no such method exists.
\end{itemize}

The mathematical result has a political moral: the 1941 Congress
made the right call.
Paradox immunity is more valuable than quota compliance because
paradoxes undermine the legitimacy of the apportionment process
in a way that occasional quota violations do not.
A state that loses a seat when the House grows has been treated
arbitrarily; a state that loses a fractional seat to rounding
has been treated consistently by a known rule.
The first is politically unacceptable; the second is mathematically
unavoidable.
